\section{Ánh xạ Kiểu dữ liệu và Ràng buộc}

\subsection{Kiểu dữ liệu trên SQL Server}
\begin{itemize}
    \item ID dùng \texttt{VARCHAR(36)} để lưu UUID do ứng dụng sinh; riêng \texttt{RoleID} dùng \texttt{INT}. UUID không thay thế kiểm soát truy cập.
    \item Tên hiển thị dùng \texttt{NVARCHAR}; email, mã khóa và mã băm dùng \texttt{VARCHAR}. Mật khẩu được băm, không lưu bản rõ.
    \item Số dư và biến động dùng \texttt{DECIMAL(18,2)}; hệ số dùng \texttt{DECIMAL(5,2)}. Thời điểm dùng \texttt{DATETIME} theo mô hình hiện tại.
    \item \texttt{SystemConfig} dùng \texttt{NVARCHAR(MAX)} nullable, kiểm tra JSON object/array bằng \texttt{ISJSON}; không sử dụng kiểu JSON riêng của hệ quản trị khác.
\end{itemize}

\subsection{Ràng buộc triển khai}
Mỗi bảng có PK; các FK dùng \texttt{ON DELETE NO ACTION}. \texttt{User.Email} duy nhất toàn hệ thống và không NULL; \texttt{ApiKey} cũng duy nhất và không NULL. Cặp \texttt{(CustomerID, MerchantID)} của CREDIT duy nhất.

\texttt{WalletAddress} và \texttt{SuiTxDigest} nullable dùng \texttt{UNIQUE INDEX ... WHERE ... IS NOT NULL}. UNIQUE thông thường trên một cột nullable của SQL Server chỉ cho phép tối đa một NULL, không phù hợp yêu cầu nhiều bản ghi chưa có định danh.

Các CHECK bảo đảm số dư không âm, hệ số dương, ngày kết thúc không trước ngày bắt đầu và miền trạng thái thống nhất. \texttt{Amount} là biến động có dấu: Earn dương, Redeem âm, Refund khác 0 và mang dấu đảo của giao dịch được hoàn. DDL chưa xác minh giao dịch gốc của Refund. \texttt{ActionType} là mã hành động mở rộng, không áp đặt danh sách đóng khi chưa có đặc tả đầy đủ.

FK tự tham chiếu của \texttt{ManagerID} kiểm tra người quản lý tồn tại; CHECK loại bỏ tự quản lý trực tiếp. Kiểm tra chu trình nhiều cấp và phạm vi Merchant thuộc tầng dịch vụ, chưa được DDL này bảo đảm.
